\section{Configuración experimental}
\label{sec:experimental}
El protocolo evaluará el aporte de D-MAB y de la cartera con controles de
igual presupuesto, siguiendo la comparación de algoritmos estocásticos
revisada en clases \parencite{vega_comparacion}. Se fijan los siguientes
criterios antes de ejecutar la campaña.

\subsection{Instancias y alcance}
Se utilizarán las 12 instancias sintéticas del repositorio oficial:
\texttt{HOSP-STD-N-01}, \texttt{-02} y \texttt{-03} para cada tamaño
de la Tabla~\ref{tab:instancias}. Se validarán esquema, elegibilidad, unidades
en minutos y huellas SHA-256. Se conservarán las restricciones y la función
objetivo de la propuesta, incluida la limpieza.

\begin{table}[ht]
 \centering
 \caption{Catálogo inicial y recursos disponibles.}
 \label{tab:instancias}
 \begin{tabular}{@{}rrrrrr@{}}
  \toprule
  Cirugías & Casos & Salas & Anestesistas & Cirujanos & Espera límite (min)\\
  \midrule
  15 & 3 & 6 & 5--6 & 7--10 & 18--55\\
  20 & 3 & 7 & 6--7 & 9--12 & 19--61\\
  25 & 3 & 8 & 7--8 & 12--15 & 21--65\\
  30 & 3 & 9 & 8--9 & 15--18 & 23--69\\
  \bottomrule
 \end{tabular}
\end{table}
Los intervalos reflejan variación entre casos y, para el límite de espera,
también entre pacientes. Este catálogo ya es conocido por el proyecto:
las nuevas semillas medirán variabilidad de búsqueda dentro de él.
La ampliación a 50, 100 y 150 cirugías se registrará como una fase separada,
con reglas de generación y admisión fijadas antes de comparar métodos;
esas instancias se reservarán de la calibración.

\subsection{Métodos y controles}
Se mantendrán dos bloques (Tabla~\ref{tab:metodos}). En A se aislarán la
selección aprendida y el detector; en B se evaluará la cartera frente
a prolongar Tabu o repetirlo. FIFO, SPT y LPT se calcularán una vez por
instancia como referencias deterministas, además de integrar la inicialización.

\begin{table}[ht]
 \centering
 \caption{Configuraciones de la campaña principal.}
 \label{tab:metodos}
 \begin{tabular}{@{}lp{0.65\textwidth}r@{}}
  \toprule
  Bloque & Método e identificador & Evaluaciones\\
  \midrule
  A & U: Tabu con selección uniforme & 3.030\\
  A & M: Tabu con D-MAB & 3.030\\
  A & Q: Tabu con UCB1 sin Page--Hinkley & 3.030\\
  B & T: Tabu uniforme prolongado & 6.060\\
  B & R: dos Tabu uniformes independientes & $3.030+3.030$\\
  B & H: cartera Tabu uniforme + Tabu con D-MAB & $3.030+3.030$\\
  \bottomrule
 \end{tabular}
\end{table}
La comparación estadística se efectuará dentro de cada bloque. La conservación
del camino uniforme de H no implica superioridad frente a T o R.

\clearpage
\subsection{Repeticiones, parámetros y parada}
Cada configuración tendrá 30 repeticiones por instancia, con semillas
80--109: se obtendrán 2.160 registros método--instancia--semilla.
U, M, Q y T compartirán las treinta soluciones iniciales de cada semilla;
H conservará el mismo inicio para sus dos caminos. R usará las semillas
$s$ y $s+1.000.000$, donde $s$ es la semilla principal.
El pareamiento se justificará por este diseño y por la instancia compartida;
los sorteos posteriores pueden divergir.

Se fijan 30 soluciones iniciales, 15 candidatos por lote y memoria de siete
soluciones. En M y Q se usará exploración $C=0{,}1$; M tendrá
margen $\delta=0{,}01$ y umbral $\lambda=1$ en Page--Hinkley.
No se ajustarán estos valores con las semillas de evaluación.
Cualquier calibración posterior utilizará datos de desarrollo y una versión
distinta del protocolo.

El presupuesto contará las valoraciones de candidatos, incluidas las
treinta iniciales de cada camino. Los duplicados sin reevaluación no gastarán
otra valoración; sus intentos y la reparación se registrarán.
La revalidación final se contabilizará aparte. Toda solución devuelta
deberá cumplir las restricciones estrictas; los fallos se conservarán en el
registro y nunca se sustituirán silenciosamente por otra semilla.

Como análisis secundario, B se comparará con 10 segundos de búsqueda:
T dispondrá de 10 s, y R y H de dos cuotas secuenciales de 5 s,
con inicialización incluida y sin transferir sobrantes.
Un plazo global de 10 s limitará ambas cuotas: el sobrepaso del primer
camino se descontará del segundo. El reloj monotónico se comprobará antes
de generar cada candidato; se registrarán el sobrepaso atómico final y las
inicializaciones incompletas. Este criterio reemplazará al de evaluaciones.
Se medirán por separado búsqueda y ejecución total, incluida la validación.

\subsection{Métricas y relevancia práctica}
La métrica primaria será $f(S)$ de \eqref{eq:objetivo}; makespan, espera
total y máxima, factibilidad, tiempo, número de evaluaciones y convergencia
se reportarán por separado. Se describirán mediana e intervalo intercuartílico,
junto con frecuencias de movimientos, reparación y reinicios del controlador.
Las curvas usarán puntos comunes del presupuesto: 10, 50 y 100\,\%.

Para comparar escalas se calculará
\begin{equation}
 g_i=100\frac{m_{c,i}-m_{a,i}}{m_{c,i}},
 \label{eq:ganancia}
\end{equation}
\noindent donde $i$ es una instancia; $m_{c,i}$ y $m_{a,i}$ son las medianas
de las treinta corridas del control y la alternativa, respectivamente,
y $g_i$ es su mejora relativa porcentual; $g_i>0$ favorece a la alternativa.
El efecto mínimo relevante se fija en \textbf{1\,\% del objetivo ponderado},
sin equipararlo a un porcentaje de minutos. La tolerancia numérica
$10^{-9}$ se aplicará a diferencias para identificar empates, conservando
los objetivos crudos. El umbral práctico y la significación se interpretarán
por separado, junto con sus intervalos de confianza.

\clearpage
\subsection{Hipótesis e inferencia estadística}
En la campaña por evaluaciones se fijan cuatro contrastes globales bilaterales:
M frente a U, H frente a T,
M frente a Q y H frente a R. El nivel será $\alpha=0{,}05$,
con corrección de Holm conjunta para las cuatro comparaciones.
La hipótesis nula de Wilcoxon será simetría de las diferencias
alrededor de cero, frente a la alternativa bilateral.
Se usará una mediana por instancia y método: \textbf{doce bloques globales},
sin tratar las 360 corridas de un método como instancias independientes.

Para inferencia se dividirá la diferencia entre ambas medianas por su máximo;
esta normalización conserva la magnitud al intercambiar métodos.
Wilcoxon usará las doce diferencias normalizadas, con permutaciones de signos
y tratamiento explícito de ceros y empates
\parencite{scipy_wilcoxon,scipy_permutation}.
Se fijará la prueba de signos como sensibilidad ante asimetría;
se reportarán ambas sin elegir la que dé menor valor $p$.
Por instancia se aplicará la misma normalización a las treinta parejas;
sus cuatro contrastes por doce instancias constituirán una familia exploratoria
de 48 pruebas, corregida con Holm.
La secundaria temporal usará las mismas semillas; H frente a T y H frente a R
formarán una familia exploratoria separada con Holm.

Como análisis contextual se aplicará Friedman a los tres métodos de cada bloque.
Por el tamaño de la muestra se usarán 100.000 permutaciones de etiquetas
dentro de cada instancia \parencite{scipy_friedman,scipy_permutation}.
Se corregirán con Holm los dos contrastes ómnibus.
Solo ante rechazo se interpretarán las tres comparaciones post hoc del bloque,
con Holm; estas no sustituirán los cuatro contrastes globales prefijados.

Se estimará la mediana de las mejoras relativas y su IC95\,\% con 5.000
remuestreos: instancias y, dentro de ellas, semillas emparejadas, conservando
juntos los métodos. Se añadirá $A_{12}$ por instancia como medida descriptiva
de superioridad marginal: más de 0,5 favorecerá a la alternativa al minimizar.
Las semillas estadísticas serán independientes y se fijarán en el protocolo.
No rechazar una hipótesis no demostrará equivalencia; treinta repeticiones
tampoco garantizan potencia suficiente para detectar una mejora del 1\,\%.

\subsection{Entorno, registros y reproducibilidad}
La referencia será macOS/arm64, Python 3.12.14, NumPy 2.3.5 y SciPy 1.18.1.
Se registrarán CPU, memoria y versiones efectivas, ejecutando secuencialmente
en la misma máquina y sin otras tareas de cálculo concurrentes.
Cada corrida guardará instancia, método, semilla, parámetros, evaluaciones,
tiempos, calidad, factibilidad y calendario final, además de trazas y eventos
del controlador.

El manifiesto incluirá huellas de datos y de todos los archivos ejecutados,
junto con la revisión del código. Se verificarán cobertura, duplicados y
calendarios antes del análisis. Los casos fallidos se describirán aparte,
con su tasa de fallos; una instancia sin resúmenes comparables no entrará
silenciosamente en una prueba.
Los contrastes exigirán treinta corridas válidas por método e instancia;
si falta alguna, se declarará incompleta la comparación correspondiente.
Se conservarán los datos sintéticos y las versiones del protocolo;
cualquier enmienda se declarará antes de su campaña.
